\documentclass[11pt, a4paper, english, twoside, notitlepage]{report}

\usepackage{iipr}
\fancyhead[RO, LE]{ }
\fancyhead[LO]{\nouppercase \leftmark}	% Heading on odd pages: Chapter name on the left
\fancyhead[RE]{\nouppercase \rightmark} % Heading on even pages: Section on the right
\fancyfoot[LE,RO]{\thepage} % Roman numbering, foot left on even, foot right on odd pages
\fancyfoot[CE]{Doble Grado en Matem\'aticas e Ingenier\'ia Inform\'atica} %Foot: Center, even pages
\fancyfoot[CO]{Universidad Complutense de Madrid (UCM)} %Foot: Center, odd pages
\renewcommand{\footrulewidth}{0.4pt} % Rule on top of foot
\renewcommand{\headrulewidth}{0.4pt} % Rule below headings
\pagestyle{fancy}

% ________Pruebas________
\usepackage{lipsum}
% _______________________

\begin{document}

\title{\huge{\textsc{Complexity analysis of polynomial algorithms}} \\
	\protect\includegraphics[scale=0.2]{ucm.pdf}}
\author{Ignacio Iker Prado Rujas \\
	Universidad Complutense de Madrid (UCM) \\
	Doble Grado en Matem\'aticas e Ingenier\'ia Inform\'atica}
\date{\today}
\maketitle
\begin{center}
	Tutors: \\
	Jos\'e F. Fernando \textit{\&} J.M. Gamboa \\
	\
\end{center}

\thispagestyle{empty}	
	
\begin{abstract}
	This work is about three different proofs of the same fact and a computational comparison between them, looking for the best one.

	Let $R$ be a real closed field and $n\geq 2$. We prove that: (1) for every finite subset $F$ of $R^n$,
	the semialgebraic set $R^n\setminus F$ is a polynomial image of $R^n$; and (2) for any
	independent linear forms $l_1,\ldots,l_r$ of $R^n$, the semialgebraic set
	$\{l_1>0,\ldots,l_r>0\}\subset R^n$ is a polynomial image of $R^n$.
	
	The key proof here is that $Q = \{x > 0, y > 0\}$ is a polynomial image of $\R^2$. This assert is proved in three different ways: a first approach using real algebraic geometry; a second and shorter one, using the composition of 3 rather simple maps; and a third one that applies topology with no computer computations.
	
	(...)
\end{abstract}

\newpage
\mbox{}
\thispagestyle{empty}    % White page
\newpage

\begin{chapter}{Introduction}
\pagenumbering{arabic}    % Numbering from 1

%\lipsum
\lipsum[14-18]
\lipsum[15]

\lipsum[16]
aslnal
\begin{section}{Hola}
\lipsum[16]
\end{section}
\lipsum[16]
\begin{Proof}
ljasjajsd $$\ceil[\Big]{\sum_{i=0}^n a_i}$$ \qed 

\end{Proof}

\end{chapter}


\end{document}